\section*{Capítulo \ref{ch:b1:membranes-transport} --- Membranas y transporte de membrana}

\begin{solution}{exo:b1:membranes-transport:1}
Una bicapa de fosfolípidos, con las colas hacia dentro y con colesterol,
forma una lámina de \qty{5}{nm}. Las proteínas están embebidas
atravesándola o adosadas a una cara, como unidades separadas. Los lípidos y
casi todas las proteínas difunden lateralmente en el plano, que es un fluido
bidimensional. Las dos monocapas difieren en composición lipídica, y las
cadenas de azúcares están solo en la cara externa.
\end{solution}

\begin{solution}{exo:b1:membranes-transport:2}
$\mathrm{O_2}$ (pequeño, apolar) $>$ etanol (pequeño, poco polar)
$>$ agua (pequeña, polar) $\gg$ glucosa (grande, polar) $\gg$
$\mathrm{Na^+}$ (con carga, hidratado). La permeabilidad cae con la
polaridad, con el tamaño y sobre todo con la carga, porque el soluto debe
disolverse en el núcleo hidrófobo.
\end{solution}

\begin{solution}{exo:b1:membranes-transport:3}
Sacarosa: $-2.48\times 0.1 = \qty{-0.25}{MPa}$. El NaCl se disocia en dos
iones, \qty{0.2}{osmol/L}: \qty{-0.50}{MPa}.
\end{solution}

\begin{solution}{exo:b1:membranes-transport:4}
$E_{Cl} = (61.5/(-1))\log(120/4) = -61.5\times 1.48 = \qty{-91}{mV}$. A
\qty{-89}{mV} el cloruro está a \qty{2}{mV} del equilibrio: se distribuye
pasivamente.
\end{solution}

\begin{solution}{exo:b1:membranes-transport:5}
Superficie celular $\num{4.7e9}\times 10^{-10}\,\mathrm{m^2} =
\qty{0.47}{m^2}$; cociente $0.92/0.47 = 1.96 \approx 2$: el lípido cubre la
\omterm{def:b1:cell-unit-of-life:cell}{célula} dos veces, luego la membrana es una bicapa.
\end{solution}

\begin{solution}{exo:b1:membranes-transport:6}
\omterm{def:b1:cell-unit-of-life:cell}{Célula}: $\Psi = -0.9 + 0.5 = \qty{-0.4}{MPa}$; disolución, \qty{-0.6}{MPa}:
el agua sale de la \omterm{def:b1:cell-unit-of-life:cell}{célula}. Equilibrio cuando $\Psi_{\text{cél}} = -0.6$:
$\Psi_p = -0.6 + 0.9 = \qty{0.3}{MPa}$; la \omterm{def:b1:cell-unit-of-life:cell}{célula} pierde algo de \omterm{def:b1:eukaryotic-cell:plantcell}{turgencia}
pero sigue turgente.
\end{solution}

\begin{solution}{exo:b1:membranes-transport:7}
Duplicar la concentración de 10 a 20 sube la velocidad solo un
\qty{15}{\%}: hay saturación. Representando $1/v$ frente a $1/c$ (o
advirtiendo que $v = 1.0$ a \qty{5}{mmol/L} y unos $1.67$ a $c$ infinita) se
obtienen $V_{\max} \approx \qty{1.7}{\micro mol/min}$ y una semisaturación
hacia \qty{3}{mmol/L}: \omterm{def:b1:membranes-transport:transporters}{difusión facilitada} por un \omterm{def:b1:membranes-transport:transporters}{transportador}.
\end{solution}

\begin{solution}{exo:b1:membranes-transport:8}
En reposo la permeabilidad de la membrana está dominada por los \omterm{def:b1:membranes-transport:transporters}{canales} de
potasio abiertos; el potasio se escapa hasta que el interior es lo bastante
negativo para retenerlo, es decir, cerca de $E_K$. Si se abrieran los
\omterm{def:b1:membranes-transport:transporters}{canales} de sodio, la permeabilidad quedaría dominada por el sodio y el
potencial se desplazaría hacia $E_{Na}$, unos \qty{+60}{mV}: el potencial de
acción.
\end{solution}

\begin{solution}{exo:b1:membranes-transport:9}
Por mol de $\mathrm{Na^+}$ que entra: $RT\ln(12/145) + F(-0.060) = -6.4 - 5.8
= \qty{-12.2}{kJ}$; dos: \qty{-24.4}{kJ}. La glucosa cuesta arriba cuesta
$RT\ln(c_{\text{int}}/c_{\text{ext}})$; la igualdad da
$\ln(c_{\text{int}}/c_{\text{ext}}) = 24\,400/2577 = 9.5$, un cociente de
unos \num{13000}.
\end{solution}

\begin{solution}{exo:b1:membranes-transport:10}
La fluidez cae (las colas se empaquetan y la bicapa se aproxima a un gel);
la difusión del oxígeno se ralentiza pero continúa; la \omterm{def:b1:membranes-transport:transporters}{bomba}, que es una
enzima, casi se detiene; las fugas continúan, de modo que a lo largo de
horas entra sodio y sale potasio y los gradientes se agotan; el potencial
cae hacia cero conforme se desvanecen los gradientes; y con sodio entrando y
la atracción osmótica de las proteínas ya no compensada, entra agua y la
\omterm{def:b1:cell-unit-of-life:cell}{célula} se hincha: las \omterm{def:b1:cell-unit-of-life:cell}{células} y los \omterm{def:b1:mammal-organization:organ}{órganos} conservados en frío se hinchan
exactamente por eso.
\end{solution}

\begin{solution}{exo:b1:membranes-transport:11}
La mezcla a \qty{37}{\celsius} y no a \qty{4}{\celsius} implica un proceso
dependiente de la temperatura: difusión en un fluido que se vuelve sólido en
frío. La ausencia de mezcla sin ATP implicaría una redistribución activa
impulsada por energía (motores, ciclado de vesículas) y no una difusión. El
resultado corregido, que la mezcla no necesita ATP, es justo lo que exige el
\omterm{def:b1:membranes-transport:membrane}{modelo del mosaico fluido}: las proteínas se desplazan por difusión térmica
en un fluido bidimensional, sin que la \omterm{def:b1:cell-unit-of-life:cell}{célula} realice trabajo.
\end{solution}

\begin{solution}{exo:b1:membranes-transport:12}
El potencial es el valor al que la fuerza eléctrica compensa la difusión del
ion más permeante, dado por Nernst; la \omterm{def:b1:cell-unit-of-life:cell}{célula} fija los gradientes (con la
\omterm{def:b1:membranes-transport:transporters}{bomba}) y las permeabilidades (con sus \omterm{def:b1:membranes-transport:transporters}{canales}), y el potencial se sigue de
ello. Una membrana de \qty{1}{\micro F/cm^2} a \qty{-70}{mV} lleva
$\qty{7e-8}{C/cm^2}$, unos $\num{4e11}$ iones por centímetro cuadrado; para
una \omterm{def:b1:cell-unit-of-life:cell}{célula} de \qty{20}{\micro m}, unos $10^5$ iones frente a los $10^{10}$
iones potasio del interior: una fracción despreciable de los iones,
desplazada una distancia despreciable, carga el condensador. Los gradientes
son la reserva; el potencial es la lectura sobre ella, y la contribución
directa de la \omterm{def:b1:membranes-transport:transporters}{bomba} (tres cargas fuera por cada dos dentro) es de unos pocos
milivoltios.
\end{solution}

\begin{solution}{pb:b1:membranes-transport:1}
\textbf{1.} $E_{Na} = 26.7\ln(145/12) = \qty{+66}{mV}$; $E_K =
26.7\ln(5/140) = \qty{-89}{mV}$.
\textbf{2.} El $\mathrm{K^+}$ (a \qty{29}{mV}) está más cerca que el
$\mathrm{Na^+}$ (a \qty{126}{mV}): la membrana es mucho más permeable al
$\mathrm{K^+}$.
\textbf{3.} $\mathrm{Na^+}$ que entra: $RT\ln(12/145) + F(-0.060) = -6.4 -
5.8 = \qty{-12.2}{kJ/mol}$. $\mathrm{K^+}$ que sale: $RT\ln(5/140) -
F(-0.060) = -8.6 + 5.8 = \qty{-2.8}{kJ/mol}$.
\textbf{4.} Coste $= 3\times 12.2 + 2\times 2.8 = \qty{42.2}{kJ}$ por
ciclo frente a \qty{50}{kJ} por ATP: rendimiento del \qty{84}{\%}.
\textbf{5.} La \omterm{def:b1:membranes-transport:transporters}{bomba} mantiene bajo el sodio citosólico; el \omterm{prop:b1:membranes-transport:secondary}{simportador} de la
cara apical usa el gradiente de sodio hacia dentro para arrastrar glucosa
desde la luz; la \omterm{def:b1:membranes-transport:transporters}{bomba} debe estar en la otra cara para que el sodio que
exporta vaya a la sangre y no de vuelta a la luz, lo que dejaría al sodio de
la luz reciclándose y el movimiento neto sería nulo.
\textbf{6.} $2\times 12.2 = \qty{24.4}{kJ/mol}$ de glucosa.
\textbf{7.} $\Delta G = RT\ln(c_C/c_L)$ (la glucosa no tiene carga).
\textbf{8.} $\ln(c_C/c_L) = 24\,400/2577 = 9.47$; $c_C/c_L =
\num{1.3e4}$.
\textbf{9.} Hasta \qty{1300}{mmol/L} en principio: el cociente nunca se
alcanza porque la glucosa sale por el \omterm{def:b1:membranes-transport:transporters}{transportador} basolateral; en la
práctica el \omterm{def:b1:eukaryotic-cell:organelle}{citosol} se mantiene entre \qtyrange{5}{20}{mmol/L}.
\textbf{10.} El \omterm{def:b1:eukaryotic-cell:organelle}{citosol} (por encima de \qty{5}{mmol/L}) está más concentrado
que la sangre; el \omterm{def:b1:membranes-transport:transporters}{transportador} deja que la glucosa difunda a favor de ese
gradiente hacia la sangre, de forma pasiva.
\textbf{11.} Al revés, el \omterm{prop:b1:membranes-transport:secondary}{simportador} bombearía glucosa desde la sangre
hacia la \omterm{def:b1:cell-unit-of-life:cell}{célula} y el \omterm{def:b1:membranes-transport:transporters}{transportador} la dejaría escapar hacia la luz: el
intestino secretaría glucosa. La polaridad de los \omterm{def:b1:membranes-transport:transporters}{transportadores} es el
sentido de la absorción.
\textbf{12.} El \omterm{prop:b1:membranes-transport:secondary}{simportador} solo capta sodio junto con glucosa; la glucosa
de la luz impulsa por tanto la absorción de sodio a través de él incluso
cuando la toxina del cólera bloquea otras vías; el cloruro sigue al sodio
por razones eléctricas y el agua sigue a la sal por \omterm{def:b1:membranes-transport:osmosis}{ósmosis}.
\textbf{13.} $\qty{4e-15}{mol/s}$ de $\mathrm{Na^+}$, es decir,
$\num{2.4e9}$ iones por segundo; a tres por ciclo, $\num{8e8}$ ciclos por
segundo.
\textbf{14.} $\num{8e8}$ ATP por segundo $= \qty{1.3e-15}{mol/s}$; a 30 ATP
por glucosa, \qty{4.4e-17}{mol/s} de glucosa, el \qty{2.2}{\%} de la glucosa
absorbida.
\textbf{15.} $\num{1.2e9}$ glucosas por segundo a 50 por \omterm{prop:b1:membranes-transport:secondary}{simportador}:
$\num{2.4e7}$ \omterm{prop:b1:membranes-transport:secondary}{simportadores}; membrana de \qty{500}{\micro m^2}: \num{48000}
por micrómetro cuadrado, es decir, uno por cada cuadrado de \qty{20}{nm}: una
membrana atestada de \omterm{def:b1:membranes-transport:transporters}{transportadores}.
\textbf{16.} Osmoles que entran: $4\times\qty{2e-15}{mol/s} =
\qty{8e-15}{osmol/s}$; a \qty{0.3}{osmol/L}, agua
$= \qty{2.7e-14}{L/s} = \qty{27}{\micro m^3/s}$. Volumen celular
$25\times 5\times 5 = \qty{625}{\micro m^3}$: duplicado en \qty{23}{s}. El
agua sale tan deprisa como entra, por la membrana basolateral, y así es como
el intestino absorbe agua.
\textbf{17.} $\qty{2.7e-14}{L/s}$ de agua son \qty{1.5e-12}{mol/s}; por glucosa ($\qty{2e-15}{mol/s}$): unas 750 moléculas de agua.
\textbf{18.} $S = 2(25\times 5 + 25\times 5 + 5\times 5) =
\qty{550}{\micro m^2} = \qty{5.5e-6}{cm^2}$; $C = \qty{5.5e-12}{F}$.
\textbf{19.} $Q = CV = 5.5\times 10^{-12}\times 0.060 =
\qty{3.3e-13}{C}$; $/\num{1.6e-19} = \num{2.1e6}$ iones.
\textbf{20.} $\num{2.1e6}/550 = 3800$ iones por micrómetro cuadrado, una carga en exceso por cada cuadrado de \qty{16}{nm} de membrana.
\textbf{21.} $\mathrm{K^+}$ en la \omterm{def:b1:cell-unit-of-life:cell}{célula}: $0.140\,\mathrm{mol/L}\times
\qty{6.25e-13}{L}\times\num{6e23} = \num{5.3e10}$ iones; la carga necesita
el \qty{0.004}{\%} de ellos: las concentraciones no cambian al cargar la
membrana.
\textbf{22.} Carga neta por ciclo, $+e$ hacia fuera: corriente
$\num{8e8}\times
\num{1.6e-19} = \qty{1.3e-10}{A}$; en un segundo la
membrana ganaría \qty{1.3e-10}{C}, es decir, $\Delta V = Q/C = \qty{24}{V}$.
No ocurre, porque cada carga bombeada hacia fuera queda compensada de
inmediato por iones que vuelven por los \omterm{def:b1:membranes-transport:transporters}{canales} (sobre todo $\mathrm{K^+}$
que se escapa menos, o $\mathrm{Na^+}$ que vuelve a entrar): la conductancia
de la membrana cortocircuita la corriente de la \omterm{def:b1:membranes-transport:transporters}{bomba}, y la \omterm{def:b1:membranes-transport:transporters}{bomba} solo añade
milivoltios.
\textbf{23.} El sodio citosólico sube en minutos; el \omterm{prop:b1:membranes-transport:secondary}{simportador} pierde su
fuerza impulsora y la absorción de glucosa se detiene; el potencial decae
hacia cero conforme se agotan los gradientes; entran sodio y agua y la
\omterm{def:b1:cell-unit-of-life:cell}{célula} se hincha.
\textbf{24.} Con la fuga de potasio bloqueada, domina la fuga de sodio: el
potencial se desplaza hacia $E_{Na}$ y se vuelve mucho menos negativo
(despolarización).
\textbf{25.} Un cociente de unos \num{13000} (\omterm{def:b1:eukaryotic-cell:organelle}{citosol} sobre luz), fijado por
dos iones sodio por glucosa a \qty{-60}{mV}; lo hacen posible la \omterm{def:b1:membranes-transport:transporters}{bomba de
sodio y potasio} (basolateral) y el \omterm{prop:b1:membranes-transport:secondary}{simportador} de sodio y glucosa (apical).
\end{solution}
